\section{Inverse of a matrix}

A square matrix $A$ represents a linear map $\mathbf{x}\mapsto A\mathbf{x}$. That map can be undone if and only if there exists a matrix $A^{-1}$, the \term{inverse matrix}, satisfying
\begin{equation}
AA^{-1}=A^{-1}A=I,
\end{equation}
where $I$ is the identity. Equivalently, in index notation,
\begin{equation}
A_{ik}(A^{-1})_{kj}=(A^{-1})_{ik}A_{kj}=\delta_{ij}.
\end{equation}
As already noted in the preliminaries, this matrix exists and is unique precisely when $\det A\neq0$. When it does, the linear system $A\mathbf{x}=\mathbf{b}$ has the unique solution
\begin{equation}
\mathbf{x}=A^{-1}\mathbf{b}.
\end{equation}

\subsection{Construction by cofactors}

To build $A^{-1}$ by hand, first form the \term{cofactor} of each entry. Deleting row $i$ and column $j$ of $A$ leaves an $(n-1)\times(n-1)$ submatrix whose determinant is the minor $M_{ij}$; the corresponding cofactor is
\begin{equation}
C_{ij}=(-1)^{i+j}M_{ij}.
\end{equation}
The \term{adjugate} (classical adjoint) is the transpose of the cofactor matrix, $(\operatorname{adj}A)_{ij}=C_{ji}$. The product identity
\begin{equation}
A\,(\operatorname{adj}A)=(\operatorname{adj}A)\,A=(\det A)\,I
\end{equation}
then yields the inverse whenever the determinant is nonzero:
\begin{equation}
\label{eq:vec-inverse-adjugate}
A^{-1}=\frac{1}{\det A}\,\operatorname{adj}A.
\end{equation}
For $n=2$ this collapses to the familiar closed form
\begin{equation}
\begin{pmatrix}a&b\\c&d\end{pmatrix}^{-1}
=\frac{1}{ad-bc}
\begin{pmatrix}d&-b\\-c&a\end{pmatrix},
\qquad ad-bc\neq0.
\end{equation}
For larger $n$ the same cofactor recipe still applies, but row reduction of the augmented matrix $[A\mid I]$ is usually faster: when the left block becomes $I$, the right block is $A^{-1}$.

If $R$ is orthogonal, $R_{ki}R_{kj}=\delta_{ij}$, then $R^{\mathsf T}R=I$ and therefore $R^{-1}=R^{\mathsf T}$. Rotation matrices are the everyday example.

\begin{example}[Inverting a $2\times2$ matrix]
Let
\begin{equation}
A=\begin{pmatrix}2&1\\1&2\end{pmatrix},\qquad
\mathbf{b}=\begin{pmatrix}1\\0\end{pmatrix}.
\end{equation}
Then $\det A=4-1=3\neq0$, so $A$ is invertible, and
\begin{equation}
A^{-1}=\frac13\begin{pmatrix}2&-1\\-1&2\end{pmatrix}.
\end{equation}
A direct multiplication confirms the definition:
\begin{equation}
AA^{-1}=\frac13\begin{pmatrix}2&1\\1&2\end{pmatrix}
\begin{pmatrix}2&-1\\-1&2\end{pmatrix}
=\frac13\begin{pmatrix}3&0\\0&3\end{pmatrix}=I.
\end{equation}
The system $A\mathbf{x}=\mathbf{b}$ is therefore solved by
\begin{equation}
\mathbf{x}=A^{-1}\mathbf{b}
=\frac13\begin{pmatrix}2&-1\\-1&2\end{pmatrix}
\begin{pmatrix}1\\0\end{pmatrix}
=\begin{pmatrix}2/3\\-1/3\end{pmatrix}.
\end{equation}
\end{example}

\begin{figure}[htb]
    \centering
    \subimport{tikz/}{inverse_map.tex}
    \caption{The inverse sends $\mathbf{b}$ back to $\mathbf{x}$. For the matrix of this example, $\mathbf{b}=\mathbf{e}_x$ and $\mathbf{x}=\tfrac23\mathbf{e}_x-\tfrac13\mathbf{e}_y$.}
    \label{fig:vec-inverse-map}
\end{figure}

\begin{example}[Cofactor inversion in three dimensions]
Consider
\begin{equation}
A=\begin{pmatrix}2&1&0\\1&2&1\\0&1&2\end{pmatrix}.
\end{equation}
Expanding along the first row gives
\begin{equation}
\det A=2(4-1)-1(2-0)=4.
\end{equation}
The cofactor matrix is
\begin{equation}
C=\begin{pmatrix}
+(4-1)&-(2-0)&+(1-0)\\
-(2-0)&+(4-0)&-(2-0)\\
+(1-0)&-(2-0)&+(4-1)
\end{pmatrix}
=\begin{pmatrix}
3&-2&1\\
-2&4&-2\\
1&-2&3
\end{pmatrix}.
\end{equation}
Here $C$ is already symmetric, so $\operatorname{adj}A=C^{\mathsf T}=C$. Equation~\eqref{eq:vec-inverse-adjugate} therefore produces
\begin{equation}
A^{-1}=\frac14\begin{pmatrix}
3&-2&1\\
-2&4&-2\\
1&-2&3
\end{pmatrix}.
\end{equation}
The $(1,1)$ entry of $A\,(\operatorname{adj}A)$ is $2\cdot3+1\cdot(-2)+0\cdot1=4=\det A$, and the off-diagonal entries cancel, in agreement with $A\,(\operatorname{adj}A)=(\det A)I$.
\end{example}
